\section{シミュレーション環境} \label{sec:simulation_environment}

本節では，本研究におけるシミュレーション環境について述べる．提案手法の有効性を検証するため，ロボットシミュレータを用いて軟弱地面上での歩行実験を行った．本節では，使用したシミュレータ，ロボットモデル，制御フレームワークおよび全身制御の構成について概説する．

\subsection{ロボットシミュレータ} \label{subsec:choreonoid}

本研究では，二足歩行ロボットの動力学シミュレーション環境として Choreonoid を用いた．Choreonoid は，剛体運動および接触力計算を含む物理シミュレーションと，ロボット制御ソフトウェアを統合的に扱うことが可能なロボットシミュレータである．物理演算エンジンには Choreonoid に標準で実装されている物理エンジンを使用し，床反力および接触拘束を考慮した運動計算を行った．

\subsection{ロボットモデル} \label{subsec:robot_model}

シミュレーションに用いたロボットモデルは，本研究で対象とする二足歩行ロボット CHIDORI のモデルである．本モデルは，※まだ

\subsection{制御フレームワーク} \label{subsec:control_framework}

歩行制御には，Choreonoid 上で動作する制御フレームワークを用いた．本フレームワークでは，歩行生成，安定化制御および全身制御が階層的に構成されている．本研究で提案する手法は，このうち重心（Center of Mass: CoM）参照を生成・調整する制御部分に実装されている．

\subsection{軟弱地面モデル}
\label{subsec:soft_ground_model}

本研究では，軟弱地面における歩行挙動を再現するため，Choreonoid 上において局所的に沈み込みを生じる床モデルを構築した．本節では，シミュレーションワールドの構成および沈み込み床の物理モデルについて述べる．

\subsubsection{シミュレーションワールドの構成}

シミュレーションワールドは，いわゆる飛び石状の足場構成とした．ロボットが歩行中に足を接地する位置にのみ床要素を配置し，それ以外の領域は歩行に影響を与えないように設定している．この構成により，ロボットが意図した接地位置に確実に足を下ろす状況を再現するとともに，床特性の違いによる影響を局所的に評価可能とした．

配置された足場のうち，一か所のみを軟弱地面として設定し，ロボットが踏み込んだ際に床が鉛直方向に沈み込むようにした．その他の足場は剛体床とし，沈み込みを生じない．このように，単一の沈み込み床を導入することで，床沈み込みが歩行制御に与える影響を明確に分離して評価することを目的とした．

\subsubsection{沈み込み床の物理モデル}

沈み込み床の力学特性は，ばね・質量・ダンパーから構成される一次元のばねマスダンパー系としてモデル化した．ロボットの足裏が床に接触すると，床要素は鉛直方向に変位し，その変位量および速度に応じた反力がロボットに作用する．この反力は，床のばね定数およびダンパー定数によって決定される．

本研究では，床の詳細な材料特性を厳密に再現することを目的とせず，沈み込み量および沈み込み速度が歩行安定性に与える影響を比較的に評価することを重視した．そのため，床パラメータは複数の組み合わせを設定し，同一条件下での比較実験を行った．

\subsubsection{床パラメータ設定}

沈み込み床に用いたばねマスダンパー系のパラメータ一覧を Table~\ref{tab:soft_ground_params} に示す．各条件では，ばね定数およびダンパー定数を変更し，沈み込みの大きさおよび応答速度の違いが歩行挙動に与える影響を評価した．なお，すべての条件において，床モデルの構成およびロボットの歩行条件は同一とした．

\begin{table}[t]
  \centering
  \caption{沈み込み床モデルに用いたパラメータ一覧}
  \label{tab:soft_ground_params}
  \begin{tabular}{c c c}
    \hline
    条件 & ばね定数 $k$ [N/m] & ダンパー定数 $d$ [Ns/m] \\
    \hline
    Case~1 & $k_1$ & $d_1$ \\
    Case~2 & $k_2$ & $d_2$ \\
    Case~3 & $k_3$ & $d_3$ \\
    Case~4 & $k_4$ & $d_4$ \\
    \hline
  \end{tabular}
\end{table}

以上のように構築した軟弱地面モデルを用いて，次節では歩行条件および評価指標について述べる．
